\label{chap:83-14}

\section{Dérivation des canaux \(90/120\) au moyen de l'inclusion des incidences hexade--octade}

La dodécade et l’alignement du
\cref{p12-83-c17-s1:sec:w12-w24-elevacion} étant fixés, soit \(H\) une hexade dans les coordonnées HMT,
soit \(\ell(H)\subset D\) son image et soit
\(O_H=\iota_D(\ell(H))\) son unique relèvement. En marquant un point
et une paire de points de l’hexade résiduelle, on obtient
\[
F_6(H)=\ell(H)\times\binom{\ell(H)}2,
\qquad
|F_6|=6\binom62=90.
\]
Si le point marqué peut parcourir l’octade :
\[
F_8(H)=O_H\times\binom{\ell(H)}2,
\qquad
|F_8|=8\binom62=120.
\]
L’inclusion \(\ell(H)\hookrightarrow O_H\) induit le morphisme d’incidence
\[
 \mathcal I_{6\to8}:F_6(H)\hookrightarrow F_8(H).
\]
En normalisant la différence par le facteur
choisi \(24\binom62\), on obtient
\[
\boxed{
\frac{90-120}{24\binom62}
=\frac{6-8}{24}
=-\frac1{12}.}
\]
L’égalité est un défaut normalisé de l’interface combinatoire
\(6\to8\). Le facteur \(24\binom62\) fait partie de la définition de cette
normalisation ; l’inclusion à elle seule détermine la différence \(-30\).
L’opérateur \(\mathcal I_{6\to8}\) a pour domaine les configurations
marquées d’incidence ; ce n’est ni l’extracteur transversal de l’état
dodécaphasique ni la famille des moments angulaires.

\section{Transport de \(-1/12\) vers les canaux \(90/120\) dans \(5\) réalisations algébriques}

Outre le défaut normalisé de l’inclusion hexade--octade, apparaissent
cinq réalisations opératorielles du même rationnel. Leurs domaines sont
précisés séparément ; l’égalité des valeurs n’identifie pas les
opérateurs sans un entrelaceur supplémentaire.

\subsection{Différence de périodicités dodécaphasiques}

Une dent de \(\Cdoce\) mesure \(30^\circ\). Alors
\[
\frac{30}{120}-\frac{30}{90}
=\frac14-\frac13
=-\frac1{12}.
\]
C’est la différence orientée entre les périodes de \(90^\circ\) et
\(120^\circ\) dans cette réalisation.

\subsection{Pseudo-inverse de l’opérateur de Gram sur la composante vivante}

Dans le saut exact \(60\to81\), deux fibres ont pour tailles \(0\) et
\(12\). L’opérateur de Gram est
\[
K_{60,81}=\operatorname{diag}(0,12).
\]
Sur la composante non nulle de son image, la pseudo-inverse vaut \(1/12\) ; avec l’orientation
négative de bord,
\[
E_{\rm surv}=-K_{60,81}^{+}
\]
produit \(-1/12\).

\subsection{2e différence entre échelles triadiques}

Définissons
\[
T_L=\sum_{n=1}^{L-1}n(L-n)=\frac{L(L^2-1)}6,
\qquad
a(L)=-\frac{T_L}{L^2}=-\frac L6+\frac1{6L}.
\]
La première soustraction d’échelle est
\[
A_1(L)=a(L)-\frac{a(3L)}3=\frac4{27L},
\]
et la seconde
\[
C(L)=LA_1(L)-\frac{3L}{3}A_1(3L)=\frac8{81}.
\]
Le résidu \(8/81\) est indépendant de \(L\). Pour obtenir
\(-1/12\), on applique la normalisation déclarée
\[
R(L)=-\frac{27}{32}C(L)=-\frac1{12}.
\]
Le calcul dérive \(8/81\) ; il ne dérive pas à lui seul le facteur \(-27/32\).

\subsection{Composition de transformations modulaires et préservation de l'état signé}

La fonction êta satisfait
\[
\frac{\eta(\tau)}{\eta(3\tau)}
=q^{-1/12}\prod_{n\ge1}(1+q^n+q^{2n})^{-1}.
\]
L’exposant provient de \(1/24-3/24=-1/12\). En prenant douze copies :
\[
t_3(q)=\left(\frac{\eta(\tau)}{\eta(3\tau)}\right)^{12}
=q^{-1}-12+54q-76q^2-243q^3+\cdots.
\]
Le coefficient \(54\) coïncide avec le défaut APP\@. L’identité rationnelle
\[
j=\frac{(t_3+27)(t_3+243)^3}{t_3^3}
\]
établit la connexion avec la branche modulaire classique.

\subsection{Défaut des projecteurs de différence cyclique}

Soit \(S\) le décalage cyclique sur \(\QQ^{12}\) et définissons
\[
 D_r=I-S^r.
\]
Le noyau de \(D_r\) est formé des vecteurs constants sur chaque orbite
de \(S^r\) ; par conséquent,
\[
 \dim\ker D_r=\gcd(12,r).
\]
Soient \(\Pi_{\ker D_3}\) et \(\Pi_{\ker D_4}\) les projecteurs rationnels sur
\(\ker D_3\) et \(\ker D_4\), et soit
\(\tau(T)=\operatorname{Tr}(T)/12\) la trace normalisée. Alors
\[
 \boxed{
 \tau(\Pi_{\ker D_3}-\Pi_{\ker D_4})
 =\frac{\operatorname{rank}\Pi_{\ker D_3}
       -\operatorname{rank}\Pi_{\ker D_4}}{12}
 =\frac{3-4}{12}
 =-\frac1{12}.}
\]
Cette réalisation exprime le rationnel comme défaut transversal des pas
trois et quatre dans le cycle à douze phases.

\begin{remark}[Domaines des réalisations]
Le défaut APP \(54\), l’exposant êta \(-1/12\), le défaut
combinatoire du relèvement hexade--octade et les
identités modulaires une fois \(t_3\) fixé sont exacts. La valeur
de la pseudo-inverse, l’extracteur normalisé et le défaut des projecteurs sont également exacts.
Chaque égalité conserve son domaine et son application correspondante.
\end{remark}


\section{Paramètre de Barbero--Immirzi : opérateur d'aire, spectre,
entrelaceur de canaux et entropie modulaire}
\label{p12-83-c14-s2:sec:area-entrelazamiento-barbero-rev11}
\index{Barbero--Immirzi!opérateur d’aire}
\index{entropie modulaire!Hamiltonien modulaire}

L'identification HMT de la fonctionnelle \(\Gamma_{6\to8}\) à la coordonnée
de Barbero--Immirzi s'effectue sur le même bord triadique qui porte
l'incidence (90/120). Le volume intérieur et son feuillet de bord possèdent la même cardinalité,
\[
 \mathcal B=(\mathbb Z/9\mathbb Z)^3,
 \quad |\mathcal B|=729,
 \qquad
 \partial\mathcal B=(\mathbb Z/27\mathbb Z)^2,
 \quad |\partial\mathcal B|=729,
\]
par le regroupement de six trits en trois coordonnées nonadiques ou deux
coordonnées d’ordre vingt-sept.

\begin{theorem}[Loi aréale de la coupe triadique]
\label{p12-83-c14-s2:thm:ley-areal-corte-triadico-rev11}
Dans le graphe cubique \(N\times N\times N\) de capacité unitaire, la coupe
minimale entre deux faces opposées a pour capacité \(N^2\). Pour \(N=9\),
\[
 \mathcal A_{\min}=81,
 \qquad S_{\rm tri}=\mathcal A_{\min}\log3=81\log3.
\]
\end{theorem}
\begin{proof}
Les \(N^2\) droites parallèles à la direction normale sont des chemins disjoints en
arêtes, de sorte que toute coupe les intercepte. Une couche transversale contient
exactement \(N^2\) arêtes et atteint la borne. Chaque liaison triadique apporte
\(\log3\) à l’entropie d’un état maximalement mélangé.
\end{proof}

Dans le tore de bord \(27\times27\), soit \(A\) le bloc canonique
\(9\times9\). Sa frontière se décompose exactement en
\begin{equation}
 \partial A=\partial A_{90}\sqcup\partial A_{120},
 \qquad |\partial A_{90}|=|\partial A_{120}|=18.
 \label{p12-83-c14-s2:eq:corte-90-120-rev11}
\end{equation}
Les liaisons horizontales représentent le canal (90), et les verticales le
canal (120). Cette décomposition transforme la paire d’incidences en une
décomposition opératorielle du bord.

L’entropie aréale \(S_{\rm tri}=81\log3\) appartient à l’état maximalement
mélangé des liaisons de la coupe. L’entropie bicouche
\(S_{\rm bi}(y)\) de la \cref{def:entropia-bicapa-rev6} appartient à la
distribution normalisée entre les deux feuillets ; les deux objets conservent
des domaines distincts.

Soit \(N=\operatorname{diag}(0,1,2)\) l’opérateur nombre d’une liaison tritique.
La réalisation aréale déclare l’échelle positive
\[
 \theta_{\rm clk}=\frac{C^*}{6}\frac{\pi}{180},
 \qquad
 a_0=\frac{1+2\cos(10^\circ)}{3}\,\theta_{\rm clk},
 \qquad \gamma_{\rm HMT}=\Gamma_{6\to8}.
\]
La formule de \(a_0\) constitue la structure de départ de cette
réalisation ; une fois cette échelle fixée, l’opérateur et son spectre se déduisent
exactement.

\begin{definition}[Opérateur d’aire HMT]
Sur le feuillet de bord
\(\mathcal H_{\partial}=\bigotimes_{e\in\partial A}\mathbb C^3\), on définit
\begin{equation}
 \widehat{\mathcal A}_{\rm HMT}
 =\gamma_{\rm HMT}a_0\sum_{e\in\partial A}N_e.
 \label{p12-83-c14-s2:eq:operador-area-hmt-rev11}
\end{equation}
Son spectre dans la coupe canonique est
\[
 \operatorname{Spec}\widehat{\mathcal A}_{\rm HMT}
 =\{n\gamma_{\rm HMT}a_0:0\leq n\leq72\},
\]
avec des multiplicités données par les coefficients de
\((1+z+z^2)^{36}\). Le prolongement par coupes compatibles produit la
famille protégée \(\mathcal A_n^{\rm prot}=n\gamma_{\rm HMT}a_0\).
\end{definition}

L’entrelaceur qui fixe la normalisation se construit dans le sous-espace des
canaux collectifs. Soient \(u_{90},u_{120}\) les vecteurs unitaires uniformes
des incidences hexadique et octadique, et soient \(v_{90},v_{120}\) les
vecteurs unitaires uniformes des deux ensembles de liaisons de
\eqref{p12-83-c14-s2:eq:corte-90-120-rev11}. L’application
\begin{equation}
 \begin{aligned}
 \mathcal J_{6\to8}:{}
 &\operatorname{span}\{u_{90},u_{120}\}
 \longrightarrow
 \operatorname{span}\{v_{90},v_{120}\},\\
 &\mathcal J_{6\to8}u_{90}=v_{90},
 \qquad
 \mathcal J_{6\to8}u_{120}=v_{120}
 \end{aligned}
 \label{p12-83-c14-s2:eq:entrelazador-area-barbero-rev11}
\end{equation}
est une isométrie entre ces sous-espaces collectifs. Si
\[
 D_{\rm inc}=90|u_{90}\rangle\langle u_{90}|
 +120|u_{120}\rangle\langle u_{120}|,
 \qquad
 D_{\rm area}=90|v_{90}\rangle\langle v_{90}|
 +120|v_{120}\rangle\langle v_{120}|,
\]
alors
\[
 \boxed{\mathcal J_{6\to8}D_{\rm inc}
 =D_{\rm area}\mathcal J_{6\to8}.}
\]
En particulier, la combinaison primitive \(12:-1\) agit
avec les mêmes coefficients avant et après \(\mathcal J_{6\to8}\). Cette
égalité fixe le dictionnaire de normalisation qui insère
\(\gamma_{\rm HMT}=\Gamma_{6\to8}\) dans la connexion
d’Ashtekar--Barbero.

L’information conservée par la réduction à la région est publiée au moyen du
Hamiltonien modulaire. Comme structure de départ de cette lecture, on déclare
\[
 \Delta_4^{\rm mod}
 =\frac{(A-C^*)\pi/180}{9\cdot90}
 \left(1-\frac7{12}\frac\pi{729}\right),
 \qquad w_{90}=90\Delta_4^{\rm mod},\quad
 w_{120}=120\Delta_4^{\rm mod},
\]
on définit
\[
 K_{\rm link}(w)=wN,
 \qquad
 \rho_{\rm link}(w)=
 \frac{\operatorname{diag}(1,e^{-w},e^{-2w})}
 {1+e^{-w}+e^{-2w}},
\]
\begin{equation}
 \rho_A=
 \bigotimes_{e\in\partial A_{90}}\rho_{\rm link}(w_{90})
 \otimes
 \bigotimes_{e\in\partial A_{120}}\rho_{\rm link}(w_{120}),
 \qquad K_A=-\log\rho_A.
 \label{p12-83-c14-s2:eq:estado-hamiltoniano-modular-barbero-rev11}
\end{equation}
Alors
\[
 K_A=\sum_{e\in\partial A}w_eN_e+c\mathbf1,
 \qquad
 S(\rho_A)=18S(w_{90})+18S(w_{120}),
\]
et l’évaluation de cette structure de départ donne
\[
 S(\rho_A)=39.54837320881808\ldots\ \text{nats}
 =57.05624190358778\ldots\ \text{bits}.
\]
La quantité \(S(\rho_A)\) est l’entropie de von Neumann du produit mixte
déclaré. Une purification globale l’identifie à une entropie
d’intrication à travers la coupe ; sans cette purification, son type canonique est
l’entropie modulaire réduite. Ainsi, \(\gamma_{\rm HMT}\) fixe le quantum
d’aire/torsion, \(S_{\rm tri}\) compte la capacité aréale, \(S_{\rm bi}\)
quantifie la répartition entre feuillets et \(S(\rho_A)\) mesure le mélange modulaire.
L’opérateur d’aire, le spectre et l’entrelaceur appartiennent à une seule construction
de bord.
